% Part: conditional-logics
% Chapter: minimal-change-semantics
% Section: true-false

\documentclass[../../../include/open-logic-section]{subfiles}

\begin{document}

\olfileid[hi]{con}{min}{tf}

\olsection{प्रतितथ्यात्मक की सत्यता और असत्यता}

यदि निकटतम $!A$-जगत सभी $!B$-जगत हों, तो प्रतितथ्यात्मक $!A \cif !B$
(गैर-रिक्ततः) सत्य है, जैसा \olref{fig:true} में चित्रित है।
\begin{figure}
\begin{center}
\begin{tikzpicture}[scale=.7]\small
  \spheresystem{5}
  \draw (0,0) node {$w$};
  \propositionintersect{0}{4}{40}{3.5}
  {\tikzset{proposition/.append={smooth,tension=1.4}}
  \proposition[shift={(1.6,-1)}]{90}{1}{30}{4}}
  \path (1.5,3) node[below] {$\formula{B}$};
  \path (3,0) node {$\formula{A}$};
\end{tikzpicture}
\caption{गैर-रिक्ततः सत्य प्रतितथ्यात्मक}
\ollabel{fig:true}
\end{center}
\end{figure}
यदि~$w$ के चारों ओर गोलक-तंत्र में कोई भी $!A$-स्वीकारी गोलक न हो, तो
प्रतितथ्यात्मक~$w$ पर भी सत्य है। उस स्थिति में वह \emph{रिक्ततः} सत्य है
(\olref{fig:vacuous} देखें)।
\begin{figure}
\begin{center}
\begin{tikzpicture}[scale=.7]\small
  \spheresystem{5}
  \draw (0,0) node {$w$};
  \proposition{0}{6.5}{40}{3.5}
  {\tikzset{proposition/.append={smooth,tension=1.4}}
  \proposition[shift={(1.2,-1)}]{90}{1}{30}{4}}
  \path (1.5,3) node[below] {$\formula{B}$};
  \spherepos{0}{7}{node {$\formula{A}$}}
\end{tikzpicture}
\caption{रिक्ततः सत्य प्रतितथ्यात्मक}
\ollabel{fig:vacuous}
\end{center}
\end{figure}

वह दो प्रकार से असत्य हो सकता है। एक प्रकार वह है जिसमें निकटतम $!A$-जगत
सभी $!B$-जगत नहीं हैं, किंतु उनमें से कुछ हैं। इस मामले में
$!A \cif \lnot !B$ भी असत्य है (\olref{fig:false} देखें)।
\begin{figure}
\begin{center}
\begin{tikzpicture}[scale=.7]\small
  \spheresystem{5}
  \draw (0,0) node {$w$};
  \propositionintersect{0}{4}{40}{3.5}
  {\tikzset{proposition/.append={smooth,tension=1.4}}
  \proposition[shift={(1.6,0)}]{90}{1}{30}{3}}
  \path (1.5,3) node[below] {$\formula{B}$};
  \path (3,0) node {$\formula{A}$};
\end{tikzpicture}
\caption{असत्य प्रतितथ्यात्मक, असत्य विपरीत}
\ollabel{fig:false}
\end{center}
\end{figure}
यदि निकटतम $!A$-जगतों का $!B$-जगतों के साथ बिल्कुल अतिव्यापन न हो, तो
$!A \cif !B$ असत्य है। किंतु इस मामले में सभी निकटतम $!A$-जगत
$\lnot !B$-जगत हैं, अतः $!A \cif \lnot !B$ सत्य है
(\olref{fig:false-opposite} देखें)।
\begin{figure}
\begin{center}
\begin{tikzpicture}[scale=.7]\small
  \spheresystem{5}
  \draw (0,0) node {$w$};
  \propositionintersect{0}{4}{40}{3.5}
  {\tikzset{proposition/.append={smooth,tension=1.4}}
  \proposition[shift={(-1,0)}]{90}{1}{30}{3}}
  \path (-1,3) node[below] {$\formula{B}$};
  \path (3,0) node {$\formula{A}$};
\end{tikzpicture}
\caption{असत्य प्रतितथ्यात्मक, सत्य विपरीत}
\ollabel{fig:false-opposite}
\end{center}
\end{figure}

कठोर सशर्त के विपरीत, प्रतितथ्यात्मक आकस्मिक हो सकते हैं।
\olref{fig:contingent} के गोलक प्रतिरूप पर विचार करें।~$u$ के निकटतम सभी
$!A$-जगत $!B$-जगत हैं, अतः $\mSat{M}{!A \cif !B}[u]$। किंतु~$v$ के
निकटतम $!A$-जगतों में कुछ $!B$-जगत नहीं हैं, अतः
$\mSat/{M}{!A \cif !B}[v]$।
\begin{figure}
\begin{center}
\begin{tikzpicture}[scale=.6]\small
  \spheresystem{7}
  \spheresystem[shift={(0:3.1)},dashed]{4}
  \propositionintersect{45}{4}{40}{4.5}
  \begin{scope}
    \clip \propositionplot{45}{4}{40}{4.5} ;
    \spherelayer[shift={(0:3.1)}]{3}
  \end{scope}
  \proposition[shift={(1.4,-.5)}]{90}{1}{35}{5}
  \path (0,0) node {$u$};
  \path (0:3.1) node {$v$};
  \spherepos{35}{9}{node {$\formula{A}$}}
  \spherepos{80}{9}{node {$\formula{B}$}}
\end{tikzpicture}
\end{center}
\caption{आकस्मिक प्रतितथ्यात्मक}
\ollabel{fig:contingent}
\end{figure}
\end{document}
